\clearpage\thispagestyle{empty}
\pdfbookmark[0]{Quellen zu den Kometenbeobachtungen von 1813}{cometa1813-sources}
\begin{center}\textbf{Quellen zu den Kometenbeobachtungen von 1813}\end{center}
\noindent Diese KI-Übersetzung folgt dem lateinischen Text in Gauß’ \textit{Werken}, Band~VI (Göttingen, 1874), S.~25--36, einschließlich des Titelblatts. Die Abhandlung wurde der Königlichen Gesellschaft der Wissenschaften am 10.~September 1813 vorgelegt und in Band~II der \textit{Commentationes societatis regiae scientiarum Gottingensis recentiores} gedruckt. \href{https://gdz.sub.uni-goettingen.de/id/PPN236007467}{Göttinger Digitalisat}.
\par\smallskip\noindent Eine eigene \href{https://www.digitale-sammlungen.de/de/view/bsb10538617?page=517}{deutsche Übersetzung aus dem Lateinischen} erschien in der \textit{Monatlichen Correspondenz zur Beförderung der Erd- und Himmels-Kunde}, Band~28, Nr.~LVIII (Gotha, Dezember 1813), S.~501--513. Der Beitrag nennt keinen Übersetzer. Auf dem Bandtitel steht F.~von Zach als Herausgeber. Das Digitalisat stammt aus der Bayerischen Staatsbibliothek.
\par\smallskip\noindent Zum Vergleich der Fachbegriffe dienten Wilhelm Olbers’ \href{https://www.e-rara.ch/zut/content/titleinfo/449756}{Abhandlung über Kometenbahnen von 1797} aus der ETH-Bibliothek und Andreas Kammerers und Mike Kretlows \href{https://fg-kometen.vdsastro.de/Publ/kometen_beobachten_2.0.pdf}{\textit{Kometen beobachten}}, Ausgabe März 2010, S.~252--256.
\par\smallskip\noindent Übersetzung und Satz wurden mit KI erstellt.
\clearpage
